\section{Late Materialization for Multi-Way Joins}
\label{sec:late_m}

The hybrid execution model described in the previous section addresses \emph{intra-operator} layout conversion overhead by separating key and payload processing within individual operators.
However, in complex analytic pipelines—particularly those involving multi-way joins—significant overhead can also arise from \emph{inter-operator} layout conversions, where data is repeatedly materialized and re-materialized as it flows between consecutive operators.

This section introduces our \emph{$N$-way late materialization} technique, which extends the hybrid execution model to eliminate redundant inter-operator conversions in multi-way join pipelines.
By deferring full tuple materialization until the final output stage, our approach avoids repeatedly reconstructing columnar output only to immediately convert it back to a row-oriented format in the next operator.

\subsection{Motivation: Inter-Operator Conversion Overhead}

Consider a common analytic pattern: an $N$-way join where the result of one join feeds into the build side of the next.
Figure~\ref{fig:nway-baseline} illustrates such a pipeline for a 3-way join of \texttt{customer}, \texttt{orders}, and \texttt{lineitem} tables (similar to TPC-H Query~3).

In the baseline execution model, each \emph{HashProbe} operator fully materializes its output—extracting matched key and payload columns from the hash table and reconstructing them in columnar form.
When this output feeds into a downstream \emph{HashBuild}, the next operator must convert these columnar vectors back into a row-oriented representation for hash table construction.
This pattern repeats at each join level, leading to $O(N)$ layout conversions for an $N$-way join.

% Even with the hybrid execution model, the inter-operator materialization remains: although payloads are stored in columnar form, the probe operator still extracts and outputs all projected columns for each matched tuple.
% When the subsequent operator is another join (or a Sort), this extracted output is immediately converted back to the hybrid format, wasting CPU cycles and memory bandwidth on intermediate materialization.

The overhead is particularly pronounced when:
\begin{itemize}[leftmargin=*]
  \item The join chain is deep (3 or more levels), amplifying redundant conversions.
  \item The projected output is wide, with many payload columns from multiple source tables.
  \item The intermediate join results are large, as occurs with low-selectivity predicates.
\end{itemize}

These observations motivate a design that defers materialization across operator boundaries, extracting payload columns only at the final stage when the output must be produced.


\subsection{Key Insight and Design Overview}

Our key insight is that in $N$-way join pipelines where join results feed into subsequent key-centric operators (another join or a sort), intermediate operators need only the \emph{key columns}—not the full payload—to perform their core processing.
Payload columns are accessed only when producing the final output.

This observation directly mirrors the insight underlying the hybrid execution model, but applied \emph{across operator boundaries} rather than within a single operator.
Whereas the hybrid model avoids materializing payload columns within an operator, $N$-way late materialization avoids materializing them between operators.

\paragraph{Design overview.}
Rather than fully materializing join output at each probe stage, our approach:
\begin{enumerate}[leftmargin=*]
  \item \textbf{Propagates only key columns} through intermediate join stages, allowing downstream operators to perform hash table construction and probing without full tuple data.
  \item \textbf{Tracks column provenance} via a \emph{columnSourceMap} that records, for each output column, where its data resides (which original table, which container).
  \item \textbf{Stores lightweight row references} instead of extracted values, enabling reconstruction of the full tuple at the final materialization point.
  \item \textbf{Defers full extraction} to the final probe operator, which retrieves all payload columns from their original sources in a single pass.
\end{enumerate}

Figure~\ref{fig:nway-late-m} illustrates this approach for the same 3-way join.
Instead of extracting and re-storing \texttt{c\_name}, \texttt{c\_address}, and other payload columns at each level, the intermediate stages propagate only the join keys (\texttt{c\_custkey}, \texttt{o\_orderkey}) along with compact references to the upstream containers.
Full tuple reconstruction occurs only at the final stage.


\subsection{Integration with the Hybrid Execution Model}

Our late materialization technique builds directly on the hybrid execution model's infrastructure.
Recall that in the hybrid model, each \emph{HybridContainer} already maintains:
\begin{itemize}[leftmargin=*]
  \item A row-oriented \emph{RowContainer} storing key columns and a compact \emph{rowRef}.
  \item Columnar \emph{RowVector}s storing payload columns.
\end{itemize}

For $N$-way late materialization, we extend the \emph{HybridContainer} with additional metadata:

\paragraph{Upstream row references.}
Each intermediate \emph{HybridContainer} stores, for every row added to its hash table:
\begin{itemize}[leftmargin=*]
  \item \texttt{upstreamBuildRowPtr}: A pointer to the matched row in the upstream join's hash table.
  \item \texttt{upstreamProbeRowId}: A compact identifier for the probe-side payload at the upstream level.
\end{itemize}

These references form a \emph{chain of provenance} linking each output row back through all upstream join levels to the original source tables.
At the final materialization point, these chains are traversed to extract columns from their sources without intermediate copying.

\paragraph{Column source map.}
A \emph{columnSourceMap} tracks, for each output column, where its data resides.
Each entry maps an output column index to a \emph{ColumnSource} descriptor containing:
\begin{itemize}[leftmargin=*]
  \item The source type: key column from a \emph{RowContainer}, payload column from columnar storage, or probe-side column.
  \item The column index within that source.
  \item A reference to the container holding the data.
\end{itemize}

This metadata is constructed incrementally as each operator in the pipeline is initialized, and is propagated through the pipeline via the execution context.
At the final probe, the complete map enables direct extraction from original sources without traversing intermediate representations.


\subsection{Execution Flow}

We now describe the execution flow through an $N$-way join pipeline with late materialization enabled.

\paragraph{Pattern detection.}
At query planning time, the execution framework detects $N$-way join patterns where a \emph{HashJoin}'s output feeds into another \emph{HashJoin}'s build side (or into a \emph{Sort}).
When such a pattern is detected, the planner marks the involved operators for late materialization and designates the \emph{final materialization point}—typically the outermost probe or the operator producing the final output.

\paragraph{Base table build (Level 0).}
The first \emph{HashBuild} in the chain operates identically to the standard hybrid model: it stores key columns in a row-oriented \emph{RowContainer} and retains payload columns in columnar form.
The \emph{columnSourceMap} is initialized to record that key columns come from this container's \emph{RowContainer} and payload columns come from its columnar storage.

\paragraph{Intermediate probe (Level 1+).}
When a \emph{HashProbe} detects that its output feeds into a downstream \emph{HashBuild} (i.e., it is not the final materialization point), it enters the \emph{late materialization output path}:
\begin{enumerate}[leftmargin=*]
  \item For each match, instead of extracting full tuples, the probe operator extracts only the row pointer from the hash table (pointing to the matched build-side row).
  \item Probe-side payload columns are stored in a \emph{ProbePayloadContainer}—a lightweight columnar container—rather than being included in the output.
  \item The output passed to the downstream \emph{HashBuild} contains only the key columns required for the next join.
  \item The \emph{columnSourceMap} is extended to include probe-side columns, mapping them to the \emph{ProbePayloadContainer}.
\end{enumerate}

\paragraph{Intermediate build (Level 1+).}
When a \emph{HashBuild} receives input from an upstream probe with late materialization enabled:
\begin{enumerate}[leftmargin=*]
  \item It extracts only the key columns needed for hash table construction from the upstream containers (using the \emph{columnSourceMap} to locate them).
  \item It stores the upstream row references (\texttt{upstreamBuildRowPtr}, \texttt{upstreamProbeRowId}) alongside each new row, maintaining the provenance chain.
  \item It updates the \emph{columnSourceMap} to reflect that key columns are now in its \emph{RowContainer}, while payload columns remain at their original sources.
\end{enumerate}

\paragraph{Final materialization.}
At the final probe operator (or the operator producing output), full tuple reconstruction occurs:
\begin{enumerate}[leftmargin=*]
  \item For each matched row in the final hash table, the operator retrieves the chain of upstream references.
  \item Using the \emph{columnSourceMap}, it determines where each output column resides.
  \item All columns are extracted from their original sources in a single pass—key columns from the appropriate \emph{RowContainer}s, payload columns from their columnar storage.
\end{enumerate}

Algorithm~\ref{algo:nway-extract} presents the final extraction procedure.


\begin{algorithm}[t]
\SetKwInOut{Input}{Input}
\Input{
  $M$: columnSourceMap (output column $\rightarrow$ source);\\
  $R$: list of matched row pointers (final level);\\
  $O$: output RowVector
}

\tcp{Phase 1: Compute row references for each level}
$L \gets$ \textsc{ComputeLevelRowRefs}($R$)\;

\tcp{Phase 2: Extract each column from its source}
\ForEach{$(col, src) \in M$}{
  \uIf{$src.type = \textsc{HybridKey}$}{
    $rows \gets L[src.level].buildRowPtrs$\;
    \textsc{ExtractFromRowContainer}($src.container$, $rows$, $src.colIdx$, $O[col]$)\;
  }
  \ElseIf{$src.type = \textsc{HybridPayload} \lor src.type = \textsc{ProbePayload}$}{
    $rowIds \gets L[src.level].rowIds$\;
    \ForEach{$i \in [0, |rowIds|)$}{
      $(cid, rid) \gets \textsc{DecodeRowId}(rowIds[i])$\;
      $O[col][i] \gets src.containers[cid][src.colIdx][rid]$\;
    }
  }
}
\caption{$N$-way late materialization: extracting columns from their original sources at the final probe stage.}
\label{algo:nway-extract}
\end{algorithm}


\subsection{Row Reference Encoding}

A key implementation detail is the encoding of row references that enable traversing the provenance chain.
Each row in an intermediate \emph{HybridContainer} stores:

\begin{itemize}[leftmargin=*]
  \item \textbf{hybridRowId}: A 64-bit value encoding the container identifier (8~bits) and local row index (56~bits), used for intra-level payload access.
  \item \textbf{upstreamBuildRowPtr}: A direct pointer (64~bits) to the matched row in the upstream hash table.
  \item \textbf{upstreamProbeRowId}: A 64-bit encoded reference to the probe-side payload at the upstream level.
\end{itemize}

This encoding enables $O(1)$ access to any column's source container without recursive traversal.
The container identifier supports up to 256 parallel drivers (sufficient for typical analytic workloads), while the 56-bit row index supports tables with up to $2^{56}$ rows.

\subsection{Handling Parallel Execution}

Modern vectorized engines execute queries using multiple parallel \emph{drivers}, each processing a disjoint subset of the input.
Late materialization must correctly handle the case where rows from different drivers are combined—for example, when the final probe matches rows that originated from different parallel instances of an upstream build.

Our design addresses this through:
\begin{itemize}[leftmargin=*]
  \item \textbf{Container identification}: Each \emph{HybridContainer} is assigned a unique identifier (the driver ID), encoded in the \emph{rowRef}.
  \item \textbf{Container merging}: Before hash table probing, containers from all drivers are merged, and each container is given visibility into all sibling containers.
  \item \textbf{Correct extraction}: During final materialization, the container identifier in each \emph{rowRef} is decoded to route extraction to the correct container.
\end{itemize}

\subsection{Discussion}

\paragraph{Comparison with traditional late materialization.}
Traditional late materialization~\cite{abadi-column-stores, monetdb-late-m} operates within a single operator, deferring column access until values are actually needed.
Our approach extends this concept across operator boundaries, but with a key difference: we maintain references not only to columnar storage (as in traditional approaches) but also to row-oriented \emph{RowContainer}s that store key columns.

This integration with the hybrid execution model is essential: in $N$-way joins, key columns must be accessible in row-oriented form for efficient hashing and comparison, yet we still wish to avoid materializing payloads.
By storing references to both key containers and payload storage, our design achieves deferred payload access while preserving row-oriented key processing.

\paragraph{Memory overhead.}
The additional memory for storing upstream references is modest.
Each intermediate row stores two additional 64-bit values (\texttt{upstreamBuildRowPtr} and \texttt{upstreamProbeRowId}), adding 16~bytes per row.
For wide analytic tables with many payload columns, this overhead is substantially smaller than the cost of materializing all payload columns at each level.

\paragraph{Applicability.}
The technique applies to any $N$-way join pattern where intermediate results feed into subsequent key-centric operators.
Common patterns include:
\begin{itemize}[leftmargin=*]
  \item Star schema queries with a large fact table probing multiple dimension table joins.
  \item Multi-table joins where smaller tables are joined first, with the result becoming the build side for larger tables.
  \item Join followed by \emph{Sort} or \emph{HashAggregation}, where the downstream operator needs only a subset of columns for its key processing.
\end{itemize}

\paragraph{Limitations.}
Late materialization is most effective when:
\begin{enumerate}[leftmargin=*]
  \item The join chain is deep (2+ levels), providing multiple opportunities to defer extraction.
  \item Payload columns are wide, making the cost of intermediate materialization high.
  \item The final selectivity is low, meaning fewer rows require full tuple reconstruction.
\end{enumerate}

For shallow joins with narrow payloads, the overhead of maintaining provenance references may outweigh the benefits.
Our implementation uses heuristics to enable late materialization only when the expected benefit exceeds the overhead.
